
\documentclass{article}
\usepackage{colortbl}
\usepackage{makecell}
\usepackage{multirow}
\usepackage{supertabular}

\begin{document}

\newcounter{utterance}

\centering \large Interaction Transcript for game `guesswhat', experiment `Level\_1', episode 7 with reka{-}flash{-}3.1{-}t0.0.
\vspace{24pt}

{ \footnotesize  \setcounter{utterance}{1}
\setlength{\tabcolsep}{0pt}
\begin{supertabular}{c@{$\;$}|p{.15\linewidth}@{}p{.15\linewidth}p{.15\linewidth}p{.15\linewidth}p{.15\linewidth}p{.15\linewidth}}
    \# & $\;$A & \multicolumn{4}{c}{Game Master} & $\;\:$B\\
    \hline

    \theutterance \stepcounter{utterance}  
    & & \multicolumn{4}{p{0.6\linewidth}}{
        \cellcolor[rgb]{0.9,0.9,0.9}{
            \makecell[{{p{\linewidth}}}]{
                \texttt{\tiny{[P1$\langle$GM]}}
                \texttt{Let's play a guessing game. Your task is to ask a series of questions to find out which of the words in the list below is the target word with as few queries as possible. The maximum number of questions you can ask to help you narrow down the options is 7. Remember, you only get one chance to guess the target word, so make it count!} \\
\\ 
\texttt{The rules of the game are: } \\
\\ 
\texttt{1. Only ask one question per turn. DO NOT ASK MORE THAN ONE QUESTION.} \\
\\ 
\texttt{2. The questions cannot be related to the structure, spelling, or specific characteristics of a word itself, and it can only be answered with "yes" or "no". } \\
\\ 
\texttt{3. Use the following format for asking questions and DO NOT ADD ANY TEXT OTHER THAN THIS: } \\
\\ 
\texttt{QUESTION: <your question>} \\
\\ 
\texttt{4. You have only one opportunity to guess the target word. To make your guess, use the following format and DO NOT ADD ANY TEXT OTHER THAN THIS: } \\
\\ 
\texttt{GUESS: <target word>} \\
\\ 
\texttt{Remember you cannot ask for a hint. Make your question brief. Let’s start!} \\
\\ 
\texttt{List of candidate words: {[}'Chrysanthemum', 'Dandelion', 'Sturgeon', 'Trout', 'Seeds', 'Edger', 'Rosemary', 'Chives'{]}} \\
            }
        }
    }
    & \\ \\

    \theutterance \stepcounter{utterance}  
    & \multicolumn{4}{p{0.6\linewidth}}{
        \cellcolor[rgb]{0.9,0.9,0.9}{
            \makecell[{{p{\linewidth}}}]{
                \texttt{\tiny{[P1$\rangle$GM]}}
                \texttt{<reasoning>} \\
\texttt{Let's see. I need to figure out which word from the list is the target using up to 7 yes/no questions. The list is {[}'Chrysanthemum', 'Dandelion', 'Sturgeon', 'Trout', 'Seeds', 'Edger', 'Rosemary', 'Chives'{]}.} \\
\\ 
\texttt{First, I should look for common categories or groupings. Let's see. Some of these are plants, others are fish, and some are tools or herbs. Let me check each word:} \\
\\ 
\texttt{{-} Chrysanthemum: a type of flower.} \\
\texttt{{-} Dandelion: a plant, also known for its seeds.} \\
\texttt{{-} Sturgeon: a type of fish.} \\
\texttt{{-} Trout: another fish.} \\
\texttt{{-} Seeds: plural, general term.} \\
\texttt{{-} Edger: a gardening tool.} \\
\texttt{{-} Rosemary: an herb.} \\
\texttt{{-} Chives: a herb.} \\
\\ 
\texttt{So possible categories: flowers, herbs, fish, gardening tools, plants (general), seeds. Maybe I can start by categorizing them. Let's see. If I can split the list into two groups, that would help. Let's think of a question that can separate these.} \\
\\ 
\texttt{First, maybe ask if the target is a type of plant. Let's see: Chrysanthemum, Dandelion, Rosemary, Chives are plants (specifically flowers or herbs). Sturgeon, Trout are animals (fish). Seeds is a plant part. Edger is a tool. So if the answer is yes, then the target is among those 5. If no, then it's either Seeds or Edger.} \\
\\ 
\texttt{Wait, but Seeds is a general term, not a specific plant. Edger is a tool. So if the first question is "Is the target a type of plant?" then possible answers:} \\
\\ 
\texttt{Yes: Chrysanthemum, Dandelion, Rosemary, Chives, Seeds.} \\
\\ 
\texttt{No: Sturgeon, Trout, Edger.} \\
\\ 
\texttt{But Seeds is a plant part, not a specific plant. Hmm. Alternatively, maybe "Is the target a living organism?" Then plants, fish, herbs, seeds would be yes. Edger is a tool, so no. But that might not split the list well. Let's see:} \\
\\ 
\texttt{Yes: Chrysanthemum, Dandelion, Sturgeon, Trout, Rosemary, Chives, Seeds.} \\
\\ 
\texttt{No: Edger.} \\
\\ 
\texttt{So that's 7 yes and 1 no. Not helpful. Maybe} \\
            }
        }
    }
    & & \\ \\

    \theutterance \stepcounter{utterance}  
    & & & \multicolumn{2}{p{0.3\linewidth}}{
        \cellcolor[rgb]{0.9,0.9,0.9}{
            \makecell[{{p{\linewidth}}}]{
                \texttt{\tiny{[GM$|$GM]}}
                \texttt{Invalid format. Guesser must use the form 'QUESTION: ' or 'GUESS: '.} \\
            }
        }
    }
    & & \\ \\

    \theutterance \stepcounter{utterance}  
    & & & \multicolumn{2}{p{0.3\linewidth}}{
        \cellcolor[rgb]{0.9,0.9,0.9}{
            \makecell[{{p{\linewidth}}}]{
                \texttt{\tiny{[GM$|$GM]}}
                \texttt{abort game} \\
            }
        }
    }
    & & \\ \\

\end{supertabular}
}

\end{document}
